\newpage
\section{Cel projektu}

Celem projektu jest stworzenie systemu wykrywania obrazów generowanych przez sztuczną inteligencję. System wykorzystuje metody uczenia maszynowego do analizy cech wizualnych i identyfikacji obrazów, które zostały wygenerowane przez modele generatywne. W kontekście współczesnych zagrożeń problem ten ma znaczenie praktyczne między innymi dla wykrywania dezinformacji, podszywania się pod inne osoby oraz obniżania wiarygodności obrazu jako nośnika informacji. W projekcie uwzględniono zarówno klasyczny etap \texttt{real vs fake}, jak i osobny etap twarzowy, przeznaczony do analizy portretów i potencjalnych deepfake'ów.

\subsection{Zakres projektu}

Zakres projektu obejmuje zaprojektowanie, implementację oraz weryfikację systemu, który klasyfikuje obrazy na prawdziwe i wygenerowane przez AI. Prace objęły przygotowanie danych, budowę modeli, analizę odporności oraz sprawdzenie zachowania modeli na danych spoza głównego rozkładu treningowego.

\begin{itemize}
    \item przygotowanie zbiorów danych i ich podział na zbiory treningowy, walidacyjny i testowy,
    \item implementacja modelu globalnego opartego o klasyfikację całego obrazu,
    \item implementacja osobnego etapu twarzowego z wykrywaniem twarzy i klasyfikacją cropów,
    \item przygotowanie benchmarku OOD dla nowszych generatorów obrazów,
    \item analiza skuteczności, odporności oraz ograniczeń modeli,
    \item przygotowanie lokalnego interfejsu demonstracyjnego do ręcznych testów i analizy wyników.
\end{itemize}

\subsection{Wymagania funkcjonalne}

\begin{itemize}
    \item wczytanie danych z ustalonej struktury katalogów \texttt{train/val/test},
    \item konfiguracja parametrów treningu i modelu z pliku konfiguracyjnego,
    \item preprocessing obrazów oraz wdrożenie augmentacji uodparniających system na popularne degradacje jakości,
    \item trening modelu oraz zapis najlepszego checkpointu,
    \item ewaluacja modelu i raport metryk \texttt{Accuracy}, \texttt{Precision}, \texttt{Recall}, \texttt{F1} i \texttt{ROC-AUC},
    \item automatyczne wydzielanie rejonu twarzy i podanie go do dedykowanego modelu twarzowego,
    \item wizualizacja przesłanek decyzyjnych modelu przy użyciu map Grad-CAM \cite{selvaraju2017gradcam}.
\end{itemize}

\subsection{Wymagania niefunkcjonalne}

\begin{itemize}
    \item przenośność --- możliwość uruchomienia na środowisku Windows/Linux,
    \item wydajność --- wsparcie dla akceleracji GPU i rozsądny czas trenowania,
    \item odtwarzalność --- zapisywanie konfiguracji, checkpointów i podsumowań eksperymentów,
    \item niezawodność --- czytelne zarządzanie wyjątkami, np. w przypadku braku wykrytej twarzy,
    \item interpretowalność --- możliwość ręcznej analizy logiki decyzyjnej modelu,
    \item bezpieczeństwo --- lokalne przetwarzanie danych bez konieczności wysyłania obrazów do usług zewnętrznych.
\end{itemize}
